\section{Introduction}\label{sec:b-intro}

Most quantum programs are not written as circuits. They are Python programs that build a circuit with a library such
as Qiskit~\cite{javadi2024qiskit}\footnote{\url{https://github.com/Qiskit/qiskit}}, deciding in loops, functions and computed indices which gates act on which qubits.
When a qubit ends in an unexpected state, the developer of such a program needs the statements that influence it, and
in a quantum program these are harder to find by reading than in a classical one. A controlled gate can change the
qubit that controls it, and a measured bit can steer operations several lines and loop iterations later. Quantum
developers report that the methods they debug with, printing, drawing circuits and reading code, do not
scale~\cite{zappin2026challenges}.

Program slicing gives a developer exactly these statements for classical programs~\cite{weiser1984program}, and
srcSlice\footnote{\url{https://github.com/srcML/srcSlice}} computes them statically, from the source and without running the program, % static-definition
for large systems on the representation of their source by srcML\footnote{\url{https://www.srcml.org}}~\cite{alomari2014srcslice,
collard2013srcml}. For a
quantum program, however, a classical slicer sees only a circuit object that every gate call uses, and the tools
that do understand circuits, such as Qiskit's DAG and its \texttt{LightCone} pass~\cite{qiskit2025}, need the program to
run first and lose the link from each gate to the line and iteration that added it. The analyses that read quantum
source code look for known patterns~\cite{paltenghi2024lintq, zhao2023qchecker}, and a qubit slicer on srcML has been
named as future work~\cite{behler2025static}.

This paper presents an approach to static slicing of quantum programs. It extends srcML and srcSlice with a quantum
layer that reads the quantum role of each call, reconstructs the circuit a program builds, over-approximating what
cannot be known without running it, slices every qubit of that circuit along the gates and measured bits that can
affect it, and shows each slice as source lines, the loop iterations that put them there and the classical statements
that compute them. We implement the approach in a tool, \tool{}, for Qiskit programs and OpenQASM~2 circuits. For
OpenQASM~2, the tool produces the srcML markup that Behler et al.\ specified~\cite{behler2024supporting} and that had so
far been written by hand, and builds on it the slicer they proposed~\cite{behler2025static}.

The paper makes four contributions.
The first is the approach, a semantic layer over srcML and srcSlice with five extensions, whose over-approximations can
only enlarge a slice, applied to Qiskit and, through an automatic front end for Behler et al.'s markup, to OpenQASM~2. The second is a way to slice loops that run on the device, with an argument that it is sound. The
third is a qubit profile, the quantum counterpart of srcSlice's slice profile. The fourth is an evaluation on
\numScaleSOneFiles{} GitHub programs, Bugs4Q~\cite{zhao2023bugs4q}, the Qiskit textbook~\cite{qiskittextbook}, Qiskit
HumanEval~\cite{vishwakarma2024qhe} and QASMBench~\cite{li2023qasmbench}. Without running a
program, \tool{} reconstructs the circuit Qiskit builds on all \numToolDevPrograms{} of our runnable programs, and its
slices match the light cone of that circuit in \numToolDevSliceEqual{} of them. It traces \numScaleSOneSliced{} GitHub
programs, takes \numScaleSOneMedianMs{}~ms per program in the median, slices every qubit of a circuit of
\numScaleQbMaxOps{} operations in \numScaleQbAllS{}~s after marking it up in \numScaleQbReadMin{} minutes, and points a developer to \numScaleReadMedianPct{} of a program's
code lines.

Section~\ref{sec:b-example} works through an example, Section~\ref{sec:b-related} reviews related work,
Section~\ref{sec:b-approach} presents the approach, Sections~\ref{sec:b-design} and~\ref{sec:b-results} the evaluation,
and the last two sections discuss its limits and conclude. The tool, the data and the scripts are in an anonymous
replication package.\TODO{anonymous package link for Paper B}
